\documentclass[11pt]{article}
\usepackage{amsmath, amssymb, bm}
\usepackage[margin=2.5cm]{geometry}
\usepackage{booktabs}

\newcommand{\hop}{\mathsf{hop\_t}}
\newcommand{\Ns}{N_s}
\newcommand{\dg}{^\dagger}
\newcommand{\cc}{\text{c.c.}}

\title{LSWT Hamiltonian: Code Convention Derivation}
\author{For internal verification}
\date{}

\begin{document}
\maketitle

%=====================================================================
\section{Basis transformation}
%=====================================================================

The code defines
\begin{equation}\label{eq:Cmat}
  \mathbf{C}
  = \frac{1}{\sqrt{2}}
    \begin{pmatrix}
      1  &  1  & 0 \\
      i  & -i  & 0 \\
      0  &  0  & \sqrt{2}
    \end{pmatrix},
  \qquad
  \mathbf{C}\dg\mathbf{C} = \mathbf{I},
  \qquad
  \mathbf{C}^T\mathbf{C}
  = \mathbf{S}
  \equiv \begin{pmatrix} 0&1&0\\1&0&0\\0&0&1 \end{pmatrix}.
\end{equation}
An important identity is
\begin{equation}\label{eq:Cstar}
  \mathbf{C}^* = \mathbf{C}\mathbf{S},
  \qquad
  \mathbf{C}^T = \mathbf{S}\,\mathbf{C}\dg.
\end{equation}

The rotated exchange matrix in the lab frame is
\begin{equation}
  \mathsf{RJ}_{\sigma\rho}
  = \mathbf{R}_\sigma^T \,\mathbf{J}_{\sigma\rho}\, \mathbf{R}_\rho,
\end{equation}
where $\mathbf{R}_\sigma$ is the rotation from the lab frame to the local
spin frame at sublattice $\sigma$.  The code computes
\begin{equation}\label{eq:hopt}
  \boxed{\hop = \mathbf{C}\dg\, \mathsf{RJ}\, \mathbf{C}.}
\end{equation}


%=====================================================================
\section{Explicit elements of $\hop$}
%=====================================================================

Direct calculation of $\hop = \mathbf{C}\dg\,\mathsf{RJ}\,\mathbf{C}$ gives:
\begin{align}
  \hop_{00} &= \tfrac{1}{2}\bigl(
    \mathsf{RJ}^{xx} + \mathsf{RJ}^{yy}
    - i(\mathsf{RJ}^{xy} - \mathsf{RJ}^{yx})\bigr),
  \label{eq:hopt00}\\
  \hop_{01} &= \tfrac{1}{2}\bigl(
    \mathsf{RJ}^{xx} - \mathsf{RJ}^{yy}
    - i(\mathsf{RJ}^{xy} + \mathsf{RJ}^{yx})\bigr),
  \label{eq:hopt01}\\
  \hop_{10} &= \tfrac{1}{2}\bigl(
    \mathsf{RJ}^{xx} - \mathsf{RJ}^{yy}
    + i(\mathsf{RJ}^{xy} + \mathsf{RJ}^{yx})\bigr),
  \label{eq:hopt10}\\
  \hop_{11} &= \tfrac{1}{2}\bigl(
    \mathsf{RJ}^{xx} + \mathsf{RJ}^{yy}
    + i(\mathsf{RJ}^{xy} - \mathsf{RJ}^{yx})\bigr),
  \label{eq:hopt11}\\[4pt]
  \hop_{02} &= \tfrac{1}{\sqrt{2}}\bigl(
    \mathsf{RJ}^{xz} - i\,\mathsf{RJ}^{yz}\bigr),
  \qquad
  \hop_{12} = \tfrac{1}{\sqrt{2}}\bigl(
    \mathsf{RJ}^{xz} + i\,\mathsf{RJ}^{yz}\bigr),
  \label{eq:hopt_col2}\\
  \hop_{20} &= \tfrac{1}{\sqrt{2}}\bigl(
    \mathsf{RJ}^{zx} + i\,\mathsf{RJ}^{zy}\bigr),
  \qquad
  \hop_{21} = \tfrac{1}{\sqrt{2}}\bigl(
    \mathsf{RJ}^{zx} - i\,\mathsf{RJ}^{zy}\bigr),
  \label{eq:hopt_row2}\\
  \hop_{22} &= \mathsf{RJ}^{zz}.
\end{align}

\paragraph{Hermiticity.}
When $\mathsf{RJ}$ is a real matrix, $\hop$ is Hermitian:
$\hop_{01} = \hop_{10}^*$, $\hop_{02} = \hop_{20}^*$, etc.


%=====================================================================
\section{The code basis: $q = \mathbf{C}\dg \bar{S}$}
%=====================================================================

The spin operator in the local frame is $\bar{S} = (\bar{S}^x, \bar{S}^y,
\bar{S}^z)^T$.  The code implicitly works in the basis
\begin{equation}\label{eq:qbasis}
  q \equiv \mathbf{C}\dg \bar{S}
  = \begin{pmatrix}
    q^0 \\ q^1 \\ q^2
  \end{pmatrix}
  = \begin{pmatrix}
    (\bar{S}^x - i\bar{S}^y)/\sqrt{2} \\[2pt]
    (\bar{S}^x + i\bar{S}^y)/\sqrt{2} \\[2pt]
    \bar{S}^z
  \end{pmatrix}
  = \begin{pmatrix}
    \tilde{S}^-/\sqrt{2} \\[2pt]
    \tilde{S}^+/\sqrt{2} \\[2pt]
    \tilde{S}^0
  \end{pmatrix}.
\end{equation}
Here $\tilde{S}^\pm = \bar{S}^x \pm i\bar{S}^y$ are the standard
raising/lowering operators.  Note the ordering:
\begin{equation}
  \boxed{
    q^0 = \frac{\tilde{S}^-}{\sqrt{2}} \;\propto\; a\dg
    \quad\text{(creation)},
    \qquad
    q^1 = \frac{\tilde{S}^+}{\sqrt{2}} \;\propto\; a
    \quad\text{(annihilation)}.
  }
\end{equation}


%=====================================================================
\section{Bilinear form in the $q$-basis}
%=====================================================================

The spin--spin interaction in the local frame is the bilinear form:
\begin{equation}
  H_{\sigma\rho} = \bar{S}_\sigma^T\; \mathsf{RJ}_{\sigma\rho}\; \bar{S}_\rho.
\end{equation}
Since $\bar{S} = \mathbf{C}\,q$ (inverting eq.~\eqref{eq:qbasis}
using $\mathbf{C}\dg\mathbf{C} = \mathbf{I}$):
\begin{equation}\label{eq:bilinear}
  H_{\sigma\rho}
  = (\mathbf{C}\,q_\sigma)^T\;\mathsf{RJ}\;(\mathbf{C}\,q_\rho)
  = q_\sigma^T\;\underbrace{\mathbf{C}^T\,\mathsf{RJ}\,\mathbf{C}}_{\displaystyle M}\;
    q_\rho.
\end{equation}
Using $\mathbf{C}^T = \mathbf{S}\,\mathbf{C}\dg$ (eq.~\eqref{eq:Cstar}):
\begin{equation}\label{eq:M}
  \boxed{
    M = \mathbf{C}^T\,\mathsf{RJ}\,\mathbf{C}
      = \mathbf{S}\;\underbrace{\mathbf{C}\dg\,\mathsf{RJ}\,\mathbf{C}}_{\hop}
      = \mathbf{S}\cdot\hop.
  }
\end{equation}
So $M_{ab} = \hop_{S(a),\,b}$, where $S$ swaps indices $0\leftrightarrow 1$:
\begin{equation}
  M_{0b} = \hop_{1b}, \qquad
  M_{1b} = \hop_{0b}, \qquad
  M_{2b} = \hop_{2b}.
\end{equation}

The bilinear is then $H = \sum_{a,b} q^a_\sigma\; M_{ab}\; q^b_\rho$,
and the coefficient of $q^a_\sigma\,q^b_\rho$ is $M_{ab}$.


%=====================================================================
\section{Holstein--Primakoff mapping}
%=====================================================================

From eq.~\eqref{eq:qbasis} and the standard HP transformation
($\tilde{S}^- \approx \sqrt{2S}\,a\dg$, $\tilde{S}^+ \approx \sqrt{2S}\,a$):
\begin{equation}
  q^0 = \frac{\tilde{S}^-}{\sqrt{2}} \approx \sqrt{S}\;a\dg,
  \qquad
  q^1 = \frac{\tilde{S}^+}{\sqrt{2}} \approx \sqrt{S}\;a,
  \qquad
  q^2 = \tilde{S}^0 \to S \;\text{(classical)}.
\end{equation}
Therefore each bilinear $q^a q^b$ maps to bosonic operators as:
\begin{center}
\begin{tabular}{llll}
\toprule
$q$-pair & Boson operator & Coupling $M_{ab}$ & $= \hop$ element \\
\midrule
$q^0\,q^0$ & $S\;a\dg a\dg$ \quad (pair creation) &
  $M_{00} = \hop_{10}$ & \texttt{hop\_t[1,0]} \\
$q^1\,q^1$ & $S\;a\,a$ \quad (pair annihilation) &
  $M_{11} = \hop_{01}$ & \texttt{hop\_t[0,1]} \\
$q^0\,q^1$ & $S\;a\dg a$ \quad (normal hopping) &
  $M_{01} = \hop_{11}$ & \texttt{hop\_t[1,1]} \\
$q^1\,q^0$ & $S\;a\,a\dg$ \quad (normal, reordered) &
  $M_{10} = \hop_{00}$ & \texttt{hop\_t[0,0]} \\[4pt]
$q^0\,q^2$ & $\sqrt{S}\;S\;a\dg$ \quad (linear, site $\sigma$) &
  $M_{02} = \hop_{12}$ & \texttt{hop\_t[1,2]} \\
$q^2\,q^0$ & $S\;\sqrt{S}\;a\dg$ \quad (linear, site $\rho$) &
  $M_{20} = \hop_{20}$ & \texttt{hop\_t[2,0]} \\
\bottomrule
\end{tabular}
\end{center}


%=====================================================================
\section{BdG Hamiltonian matrix}
%=====================================================================

The BdG structure (PDF eq.~40, 45):
\begin{equation}
  \hat{\mathsf{H}}_\mathbf{k}
  = \begin{pmatrix}
    \mathsf{A}_\mathbf{k}  & \mathsf{B}_\mathbf{k} \\[2pt]
    \mathsf{B}\dg_\mathbf{k} & \mathsf{A}^*_{-\mathbf{k}}
  \end{pmatrix},
  \qquad
  \Psi\dg_\mathbf{k} = (a\dg_{\mathbf{k},1},\ldots,a\dg_{\mathbf{k},\Ns},\;
  a_{-\mathbf{k},1},\ldots,a_{-\mathbf{k},\Ns}).
\end{equation}

\begin{itemize}
\item Upper-right $[\sigma,\Ns+\rho]$ = $\mathsf{B}_\mathbf{k}$:
  couples $a\dg_{\mathbf{k},\sigma}\,a\dg_{-\mathbf{k},\rho}$
  = pair creation $\Rightarrow$ coefficient $M_{00} = \hop_{10}$.
\item Lower-left $[\Ns+\sigma,\rho]$ = $\mathsf{B}\dg_\mathbf{k}$:
  Hermitian conjugate of $\mathsf{B}_\mathbf{k}$.
\end{itemize}

\paragraph{Code assignment.}
The code defines \texttt{tpp = $\sqrt{S_\sigma S_\rho}\;\hop_{10}$}
and assigns:
\begin{align}
  \mathsf{H}[\sigma,\,\Ns{+}\rho]
  &\mathrel{+}= \texttt{tpp}\cdot e^{-i\mathbf{k}\cdot\bm{\delta}}
  \qquad\text{(B block, upper-right)},\\
  \mathsf{H}[\Ns{+}\rho,\,\sigma]
  &\mathrel{+}= (\texttt{tpp}\cdot e^{-i\mathbf{k}\cdot\bm{\delta}})^*
  \qquad\text{(B$\dg$ block, lower-left)}.
\end{align}
This is \textbf{correct}: the pair creation coupling $M_{00} = \hop_{10}$
appears in the B block.

\paragraph{Normal hopping.}
The code defines \texttt{tpm = $\sqrt{S_\sigma S_\rho}\;\hop_{11}$}:
\begin{align}
  \mathsf{H}[\sigma,\rho]
  &\mathrel{+}= \texttt{tpm}\cdot e^{-i\mathbf{k}\cdot\bm\delta}
  \qquad\text{(A block)},\\
  \mathsf{H}[\rho,\sigma]
  &\mathrel{+}= (\texttt{tpm}\cdot e^{-i\mathbf{k}\cdot\bm\delta})^*
  \qquad\text{(Hermitian conjugate)}.
\end{align}
This is \textbf{correct}: the normal hopping coupling $M_{01} = \hop_{11}$
appears in the A block.


%=====================================================================
\section{Linear terms}
%=====================================================================

The linear term is the coefficient of $a\dg_\sigma$ in the Hamiltonian,
which must vanish at classical equilibrium.

\paragraph{From exchange couplings.}
For site $\sigma$ from the $q^0_\sigma\,q^2_\rho$ coupling
(with $q^2_\rho \to S_\rho$ classically):
\begin{equation}
  \text{coupling } M_{02} = \hop_{12},
  \qquad
  q^0_\sigma \approx \sqrt{S_\sigma}\;a\dg_\sigma.
\end{equation}
For site $\rho$ from $q^2_\sigma\,q^0_\rho$:
\begin{equation}
  \text{coupling } M_{20} = \hop_{20},
  \qquad
  q^0_\rho \approx \sqrt{S_\rho}\;a\dg_\rho.
\end{equation}

\paragraph{Code assignment.}
\begin{align}
  f_\sigma
  &\mathrel{+}= \sqrt{2}\; S_\rho\; \hop_{12}^{(\sigma\rho)},\\
  f_\rho
  &\mathrel{+}= \sqrt{2}\; S_\sigma\; \hop_{20}^{(\sigma\rho)}.
\end{align}
The $\hop$ indices \texttt{[1,2]} and \texttt{[2,0]} match $M_{02}$
and $M_{20}$ respectively.
The $\sqrt{2}$ prefactor arises from the relation
$\mathsf{RJ}^{xz} + i\,\mathsf{RJ}^{yz} = \sqrt{2}\;\hop_{12}$
(eq.~\eqref{eq:hopt_col2}), converting between the Cartesian and
$q$-basis normalizations.
At classical equilibrium, $f_\sigma = 0$ for all sublattices.


%=====================================================================
\section{Comparison with the PDF note and legacy code}
%=====================================================================

\paragraph{PDF note (eq.~28).}
The PDF defines $\tilde{J}^{\alpha\beta}$ with $\alpha,\beta \in \{+,-,0\}$
via $(C\dg\,\mathsf{RJ}\,C)_{\alpha\beta}$ (eq.~27), with indices
ordered as $(+,-,0)$ in the matrix display.  However, the $q$-basis
(eq.~\eqref{eq:qbasis}) has ordering $(-,+,0)$ for its indices $(0,1,2)$.
Comparing with the explicit formulas in eq.~28:
\begin{equation}
  \hop_{00} = \tilde{J}^{-+}, \quad
  \hop_{01} = \tilde{J}^{--}, \quad
  \hop_{10} = \tilde{J}^{++}, \quad
  \hop_{11} = \tilde{J}^{+-}, \quad
  \hop_{22} = \tilde{J}^{00}.
\end{equation}
The PDF bilinear (eq.~23) uses $\tilde{J}^{\bar\alpha,\beta}$ with
$\bar{+}=-$, $\bar{-}=+$.  This $\bar\alpha$ convention compensates
for the $(+,-,0)$ labeling vs the $q$-basis $(-, +, 0)$ ordering:
\begin{equation}
  \text{coupling for }q^a q^b
  = \tilde{J}^{\overline{\hat\alpha(a)},\,\hat\beta(b)}
  = M_{ab} = \hop_{S(a),b},
\end{equation}
where $\hat\alpha(0) = -$, $\hat\alpha(1) = +$, and $\bar{-} = +$, so
$\overline{\hat\alpha(0)} = +$ selects row~0 of $\tilde{J}$, which is
row~1 of $\hop$ (via the S-swap).  All approaches give the same result.

\paragraph{Legacy code.}
The legacy code uses the same \texttt{tpp = hop\_t[1,0]} but places
it in the \textbf{lower-left} block (B$\dg$ position) with its
conjugate in the upper-right:
\begin{align}
  \text{Legacy:}\quad
  \mathsf{H}[\Ns{+}\sigma,\,\rho] &\mathrel{+}= \texttt{tpp}\cdot e^{-i\mathbf{k}\cdot\bm\delta}
  \quad\text{(lower-left)},\\
  \mathsf{H}[\rho,\,\Ns{+}\sigma] &\mathrel{+}= \texttt{tpp}^*\cdot e^{+i\mathbf{k}\cdot\bm\delta}
  \quad\text{(upper-right)}.
\end{align}
This swaps B and B$\dg$: the upper-right (pair creation) block
receives $\hop_{10}^* = \hop_{01}$ instead of $\hop_{10}$.
For the stripe phase where $\mathsf{RJ}$ is real symmetric and
$\hop_{10} = \hop_{01}^*$ is real, B$\,{=}\,$B$\dg$ and the swap has
no effect.  For non-collinear phases (e.g., NBCP) where
$\hop_{10} \ne \hop_{10}^*$, the legacy code produces incorrect results.


%=====================================================================
\section{Summary}
%=====================================================================

\begin{center}
\begin{tabular}{llll}
\toprule
Physical quantity & Coupling & $\hop$ index & Code variable \\
\midrule
Pair creation (B block, $a\dg a\dg$)
  & $M_{00}$ & \texttt{[1,0]} & \texttt{tpp} \\
Normal hopping (A block, $a\dg a$)
  & $M_{01}$ & \texttt{[1,1]} & \texttt{tpm} \\
Diagonal ($S^z S^z$)
  & $M_{22}$ & \texttt{[2,2]} & \texttt{t00} \\
Linear (site $\sigma$, $a\dg_\sigma$)
  & $M_{02}$ & \texttt{[1,2]} & --- \\
Linear (site $\rho$, $a\dg_\rho$)
  & $M_{20}$ & \texttt{[2,0]} & --- \\
\bottomrule
\end{tabular}
\end{center}

All code indices are correct.  The relationship $M = \mathbf{S}\cdot\hop$
(eq.~\eqref{eq:M}) provides the systematic mapping between the physical
coupling matrix and the code's \texttt{hop\_t} elements.

\end{document}
